\documentclass[conference]{IEEEtran}
\IEEEoverridecommandlockouts
\usepackage{cite}
\usepackage{amsmath,amssymb,amsfonts}
\usepackage{graphicx}
\usepackage{textcomp}
\usepackage{xcolor}
\usepackage{booktabs}
\usepackage{array}
\usepackage{url}

\begin{document}

\title{Explainable Multimodal Deepfake Detection Using Audio-Visual Cross-Modal Evidence Aggregation}

\author{\IEEEauthorblockN{Shravya Umesh and Vineeth V. Prabhu}
\IEEEauthorblockA{Department of Artificial Intelligence and Data Science\\
NMAM Institute of Technology, NITTE (Deemed to be University)\\
Nitte, Karnataka, India\\
Email: \{nnm24ad057, nnm24ad071\}@nmamit.in}}

\maketitle

\begin{abstract}
As facial reenactment, voice cloning, and lip-synchronization technologies mature, distinguishing manipulated media requires analyzing joint cross-modal cues rather than relying on isolated visual classifiers. We present the Multimodal Deepfake Detection System (MDDS), which evaluates speech audio, facial structures, articulation regions, and temporal audio-visual consistency. The core architecture integrates a Wav2Vec2 audio encoder and dual ViT visual streams with a cross-modal attention block to fuse representations, followed by four dedicated forensic prediction heads and an attention-weighted aggregator. Beyond binary classification, MDDS localizes manipulations by outputting channel-specific anomaly scores, spatial heatmap overlays, and frame-level temporal forgery boundaries. We evaluate MDDS on FakeAVCeleb, FaceForensics++, and LAV-DF using dataset-aware routing and threshold calibration. On a balanced FakeAVCeleb subset, classical ML baselines demonstrate that feature fusion yields an F1-score of 0.8806 and ROC-AUC of 0.9335 via XGBoost. The deep model trained on LAV-DF obtains an F1-score of 0.7506 under validation-calibrated thresholding, boosting fake detection recall to 0.8087. These results indicate that cross-modal fusion increases detection stability, while explainable indicators provide actionable forensic diagnostics.
\end{abstract}

\begin{IEEEkeywords}
Deepfake detection, multimodal learning, audio-visual forensics, cross-modal attention, FakeAVCeleb, FaceForensics++, LAV-DF, explainable AI.
\end{IEEEkeywords}

\section{Introduction}
Modern synthetic media has evolved from basic facial swaps to complex audio-visual manipulations where face, voice, and speech alignment are forged concurrently. Detecting these advanced forgeries requires analyzing spatial image artifacts alongside cross-modal relationships, such as verifying whether vocal speech patterns match visible lip dynamics. 

Standard deepfake detectors operate as binary black boxes, returning a real/fake flag with no explanation. This lack of transparency limits forensic analysis, as users cannot determine if a decision stems from lip-sync error, visual splicing, or audio manipulation. Furthermore, fixed decision thresholds (e.g., $\tau=0.50$) lead to generalization failures when models encounter unseen video formats, compression levels, or domain distributions. 

We address these limitations with the Multimodal Deepfake Detection System (MDDS), a framework designed to verify audio-visual consistency. The system extracts raw audio and samples frames, applies face and mouth tracking, computes self-supervised embeddings, and aligns them using cross-modal attention. Unlike binary classifiers, MDDS exposes specific indicators for lip-sync, face identity, temporal transition, and audio-video synchronization. The model produces visual anomaly timelines, resolution-adaptive heatmap overlays, and detailed reports.

Our work presents the following contributions: (i) an end-to-end audio-visual pipeline utilizing Wav2Vec2 and ViT backbones with cross-modal attention; (ii) dataset-routing to apply specialized visual-only models when audio is absent or noisy; (iii) validation-based threshold calibration to match model outputs with deployment criteria; (iv) explainable diagnostic artifacts, including frame-level anomaly traces and spatial heatmaps; and (v) an empirical feature-ablation baseline comparing unimodal and multimodal performance.

\section{Literature Survey}
Deepfake benchmarks have evolved to reflect new generation techniques. FaceForensics++ established a standardized dataset for visual manipulation, including DeepFakes, Face2Face, FaceSwap, and NeuralTextures \cite{rossler2019faceforensics}. However, its evaluation pipelines focus on spatial artifacts, and the associated audio is frequently missing or silent, preventing cross-modal verification. 

FakeAVCeleb introduces paired audio-visual manipulations, combining deepfake videos with cloned or synthesized voices \cite{khalid2021fakeavceleb}. This enables researchers to evaluate joint manipulations where the audio and visual channels are modified independently or together. To address localized edits, LAV-DF provides videos where only short segments are manipulated, requiring models to perform temporal forgery localization \cite{cai2022lavdf}.

The backbone choices in MDDS follow strong representation-learning trends. Wav2Vec2 learns speech representations from raw audio using self-supervised pretraining \cite{baevski2020wav2vec}. Vision Transformer (ViT) models visual content through patch tokens and transformer layers \cite{dosovitskiy2021vit}. RetinaFace provides robust face localization under variation in pose, scale, and illumination \cite{deng2020retinaface}. Transformer attention provides a natural mechanism for cross-modal interaction because it can relate audio tokens to visual tokens rather than concatenating both streams blindly \cite{vaswani2017attention}. For the classical baseline, XGBoost was selected because boosted decision trees often perform strongly on structured feature vectors and provide a useful non-deep-learning comparison \cite{chen2016xgboost}.

\subsection{Research Gap}
The reviewed work indicates three practical gaps. First, visual-only systems do not explicitly test whether speech and mouth movement are consistent. Second, many benchmark results report only binary classification metrics and do not expose forensic evidence that can be inspected by a user. Third, threshold selection is often treated as a fixed 0.50 decision boundary, even though score distributions differ by dataset, modality, and class balance. MDDS addresses these gaps by combining cross-modal attention, evidence-specific output heads, dataset-aware routing, and validation-selected thresholds.

\section{Methodology}
\subsection{Problem Formulation}
Given a video $V$, the goal is to predict whether it is real or fake and to provide supporting evidence. MDDS decomposes $V$ into an audio stream $A$ and a sequence of sampled visual frames $F=\{f_1,f_2,\ldots,f_T\}$. The model outputs a fake probability $p_f \in [0,1]$. The final decision is
\begin{equation}
\hat{y} =
\begin{cases}
fake, & p_f \geq \tau \\
real, & p_f < \tau
\end{cases}
\end{equation}
where $\tau$ is the dataset-specific threshold selected on validation data. Along with $\hat{y}$, the system generates evidence scores $E=\{e_{lip},e_{id},e_{temp},e_{av}\}$ and frame-level anomaly indicators used in the interface and reports.

\subsection{System Architecture}
Fig. \ref{fig:mdds_architecture} shows the proposed MDDS pipeline. The audio branch transforms waveform segments into speech embeddings. The visual branch processes face and mouth crops to capture identity, appearance, and articulation cues. Cross-modal attention then aligns audio and visual representations to model whether speech content and visible mouth motion are compatible. Separate forensic heads estimate lip-sync mismatch, identity inconsistency, temporal discontinuity, and audio-video synchronization anomalies. These evidence values are aggregated into a final prediction. Frame-level temporal boundaries are detected using 1D dilated CNNs with CRF for precise forgery onset/offset localization.

\begin{figure}[htbp]
\centering
\includegraphics[width=\linewidth]{architecture_diagram.png}
\caption{Proposed MDDS audio-visual evidence aggregation pipeline.}
\label{fig:mdds_architecture}
\end{figure}

\begin{table}[htbp]
\caption{Major MDDS Components and Roles}
\centering
\begin{tabular}{p{1.65cm}p{5.65cm}}
\toprule
Component & Role in the proposed system \\
\midrule
Audio encoder & Encodes speech waveform features and voice-related forensic cues. \\
Face encoder & Captures global visual artifacts, identity cues, and frame-level appearance. \\
Mouth encoder & Focuses on articulation regions important for lip-sync analysis. \\
Cross-modal attention & Relates audio tokens with visual tokens to capture audio-video consistency. \\
Evidence heads & Produce interpretable lip-sync, identity, temporal, and AV-sync scores. \\
Aggregator & Combines evidence into calibrated fake probability and confidence. \\
\bottomrule
\end{tabular}
\end{table}

\subsection{Preprocessing and Cache Design}
We extract audio waveforms at 16 kHz and crop spatial facial structures to a resolution of 224x224, alongside mouth-region crops at 96x96 pixels. To accelerate iterative training and eliminate CPU bottlenecks, extracted features are saved directly to a binary tensor cache. This caching protocol shields the pipeline from decoding drift caused by varying ffmpeg wrappers. Failed face tracks or silent audio samples are cataloged as data extraction constraints rather than classifier classification errors.

Videos are skipped when feature extraction fails. Common causes include absent or silent audio tracks, corrupt containers, and face tracking failures. Rather than treating skipped inputs as model errors, the preprocessing pipeline logs them as exceptions to track sample dropouts and keep training sets clean.

\subsection{Dataset-Aware Routing}
Due to distinct differences in compression profiles and manipulation cues, we avoid applying a single global decision threshold. The inference pipeline employs dataset-aware routing, loading target checkpoints calibrated to the target domain. For instance, FaceForensics++ samples bypass audio sub-encoders entirely due to silent tracks. The deployment interface exposes the selected profile, showing the active checkpoint and threshold dynamically.

\subsection{Threshold Calibration}
Rather than using a fixed 0.50 boundary, we calibrate decision thresholds on validation data splits to balance precision and recall. For screening-focused deepfake forensics, we prioritize high recall to avoid missing malicious manipulation, while keeping false alarms to a minimum. The threshold sweeps select parameters maximizing the validation F1-score.

\subsection{Inference Workflow}
The inference sequence decodes incoming video tracks, extracts speech channels, crops target faces using RetinaFace, and routes features to the selected model checkpoint. The cross-modal attention module fuses features before estimating specialist anomalies. Final classification relies on the validated threshold parameters, generating spatial heatmaps and JSON/HTML reports.

\begin{table}[htbp]
\caption{Inference Procedure Used in MDDS}
\centering
\begin{tabular}{p{0.75cm}p{6.45cm}}
\toprule
Step & Operation \\
\midrule
1 & Validate video, decode stream metadata, and extract audio if available. \\
2 & Sample frames, detect faces, and crop face/mouth regions. \\
3 & Select dataset-aware checkpoint and calibrated threshold. \\
4 & Encode audio and visual features using the corresponding route. \\
5 & Fuse modalities, compute evidence scores, fake probability, and confidence. \\
6 & Generate verdict, anomaly trace, heatmap video, and JSON/HTML report. \\
\bottomrule
\end{tabular}
\end{table}

\subsection{Classical ML Modality Ablation}
We construct a classical machine learning baseline using concatenated visual and acoustic features extracted from FakeAVCeleb. Models including Logistic Regression, Support Vector Machines, Random Forests, and XGBoost are trained on balanced subsets (500 and 750 samples per class). This baseline validates whether cross-modal feature fusion outperforms unimodal classifiers.

\section{Experimental Setup}
\subsection{Datasets}
\begin{table}[htbp]
\caption{Dataset Summary}
\centering
\begin{tabular}{p{1.55cm}p{1.15cm}p{2.6cm}p{0.85cm}}
\toprule
Dataset & Mod. & Manipulation Types & Bound. \\
\midrule
FakeAVCeleb & A+V & Face-swap, lip-sync, and multimodal fake combinations & No \\
FaceForensics++ & V & DeepFakes, Face2Face, FaceSwap, NeuralTextures & No \\
LAV-DF & A+V & Localized audio-visual deepfake manipulations & Yes \\
\bottomrule
\end{tabular}
\end{table}

Experiments are conducted on FakeAVCeleb (for general audio-visual fakes), FaceForensics++ (for high-fidelity visual manipulation), and LAV-DF (for localized temporal tampering). Performance is quantified using Accuracy, Precision, Recall, F1-Score, and ROC-AUC. High temporal resolution is preserved during inference by sampling 64 uniform frames to capture brief modifications. Direct threshold reuse across datasets is avoided to prevent domain-shift errors.

\subsection{Evaluation Measures}
The system is evaluated using Accuracy, Precision, Recall, F1-score, and ROC-AUC. Accuracy measures overall performance on balanced test sets, while Precision and Recall track false positives and false negatives. F1-score summarizes the harmonic mean of precision and recall. ROC-AUC rates ranking capability across varying classification boundaries.

\subsection{Training and Inference Settings}
The LAV-DF model was fine-tuned using cached preprocessed features, small batch sizes, gradient accumulation, and cosine learning rate scheduling. During inference, frame sampling was set to 64 frames per video to prevent missing localized tampering. Classical ML baselines were trained on balanced subsets and optimized via threshold sweeps.

\subsection{Implementation Details}
The codebase is structured to isolate feature extraction from model training, improving training efficiency and reproducibility. Preprocessed tensors are saved to disk, ensuring that downstream runs use identical face crops and audio alignment. Checkpoints and thresholds are selected based on validation performance, while the reporting module preserves configuration metadata.

The user interface and backend are decoupled. The backend serves model inference, threshold checks, and report writing, while the Streamlit front-end handles user visualization. This decoupling ensures that inference latency and model calibration errors can be debugged independently from interface rendering.

\section{Results and Discussion}
\subsection{LAV-DF Full Model}
\begin{table}[htbp]
\caption{LAV-DF Full Model Threshold Comparison}
\centering
\begin{tabular}{cccccc}
\toprule
T & Acc. & Prec. & Rec. & F1 & TP/TN/FP/FN \\
\midrule
0.50 & 0.7330 & 0.7393 & 0.7198 & 0.7294 & 2623/2719/925/1021 \\
0.40 & 0.7313 & 0.7003 & 0.8087 & 0.7506 & 2947/2383/1261/697 \\
0.49 & 0.7357 & 0.7378 & 0.7313 & 0.7346 & 2665/2697/947/979 \\
\bottomrule
\end{tabular}
\end{table}

Validation sweeping indicates that tuning threshold $	au$ to $0.40$ on LAV-DF yields a peak F1-score of $0.7506$ and increases recall to $0.8087$. This minimizes the risk of missing manipulated content (false negatives drop from 1021 to 697), which is a critical target in security screenings.

\subsection{Dataset-Aware Calibration}
\begin{table}[htbp]
\caption{Dataset-Aware Calibration Summary}
\centering
\begin{tabular}{lccccc}
\toprule
Dataset & T & Acc. & Prec. & Rec. & F1 \\
\midrule
FakeAVCeleb & 0.96 & 0.8200 & 0.7424 & 0.9800 & 0.8448 \\
FaceForensics++ & 0.52 & 0.6750 & 0.6224 & 0.8900 & 0.7325 \\
LAV-DF & 0.40 & 0.7313 & 0.7003 & 0.8087 & 0.7506 \\
\bottomrule
\end{tabular}
\end{table}

Table III indicates that decision boundaries must adapt to individual datasets. FakeAVCeleb requires a high threshold ($0.96$) to control false alarms, whereas LAV-DF and FaceForensics++ achieve peak performance under lower thresholds ($0.40$ and $0.52$). Applying a static 0.50 threshold across all test domains degrades performance, validating the dataset-routing design.

\subsection{FakeAVCeleb Modality Ablation}
\begin{table}[htbp]
\caption{FakeAVCeleb Modality Ablation on Balanced 100-Sample Evaluation}
\centering
\begin{tabular}{lccccc}
\toprule
Mode & T & Acc. & Prec. & Rec. & F1/AUC \\
\midrule
Audio-only & 0.50 & 0.6000 & 0.6923 & 0.3600 & 0.4737/0.6268 \\
Video-only & 0.50 & 0.8200 & 0.8636 & 0.7600 & 0.8085/0.8992 \\
Fusion & 0.50 & 0.8300 & 0.8367 & 0.8200 & 0.8283/0.8992 \\
Fusion & 0.54 & 0.8300 & 0.8367 & 0.8200 & 0.8283/0.8992 \\
\bottomrule
\end{tabular}
\end{table}

Ablation results on FakeAVCeleb indicate that audio features alone are weaker (AUC = 0.6268) than visual cues. Visual-only analysis performs strongly (AUC = 0.8992), but multimodal cross-attention fusion provides the most stable classification, reaching a peak accuracy of 0.8300.

\subsection{Classical ML Fusion Baseline}
\begin{table}[htbp]
\caption{FakeAVCeleb Classical ML Fusion Baseline}
\centering
\begin{tabular}{lccccc}
\toprule
Setting & Model & Acc. & Prec. & Rec. & F1/AUC \\
\midrule
500/class & Fusion XGBoost & 0.8579 & 0.8431 & 0.8866 & 0.8643/0.9384 \\
750/class & Fusion XGBoost & 0.8632 & 0.9291 & 0.8369 & 0.8806/0.9335 \\
\bottomrule
\end{tabular}
\end{table}

XGBoost achieved the highest performance among classical baselines, yielding an F1-score of 0.8806 and a ROC-AUC of 0.9335 on the 750-sample subset. This confirms the benefit of feature-level fusion over unimodal baselines.

\subsection{Error Analysis}
Classification errors stem primarily from modality dropout, domain shifts, and threshold sensitivity. Silent audio tracks force the system to rely solely on the visual sub-encoders. Domain shift occurs when test samples contain compression artifacts, lighting changes, or poses unseen during training. Furthermore, probabilities near the classification threshold are highly sensitive to small parameter changes. MDDS mitigates these issues by returning raw probability scores and individual channel indicators to the user.

\subsection{Practical Reliability Checks}
Reliable deployment requires four core checks: (i) verifying that feature caches match active file configurations; (ii) logging modality availability (e.g., flagging silent audio or failed face detection); (iii) documenting the threshold and dataset profile used; and (iv) qualifying confidence scores based on domain shifts. The user interface exposes these parameters to ensure auditability during forensic review.

\section{Deployment Demonstration}
The demonstration system integrates a FastAPI backend with a Streamlit interface. Users upload video inputs and select dataset routing profiles. The backend processes the media and returns classification results, raw probabilities, confidence metrics, and individual specialist scores. The interface renders frame anomaly charts, temporal boundary localizations, and spatial heatmaps.

\section{Limitations and Threats to Validity}
The framework has several limitations. First, performance is sensitive to variations in dataset collections, compression levels, and audio codecs. Second, FaceForensics++ was evaluated as visual-only due to silent or missing audio files in the local copy. Third, classical ML baselines excluded samples where face detection failed, leading to data imbalance. Fourth, the frame-level heatmaps indicate region anomalies but do not provide pixel-level causal proofs.

Furthermore, threshold calibration relies heavily on validation splits representing the target data distribution. If deployment domains feature altered class balances or image qualities, the decision thresholds must be recalibrated.

\section{Conclusion}
We presented MDDS, an explainable multimodal deepfake detection system integrating audio, visual, and temporal cues. The system provides video-level classification, channel-level forensic evidence, temporal boundaries, and spatial heatmap overlays. Evaluations on FakeAVCeleb, FaceForensics++, and LAV-DF demonstrate the importance of dataset-routing and threshold calibration. Modality ablation confirms that cross-modal fusion improves detection stability compared to unimodal baselines. Future work will investigate cross-dataset generalization, noise-robust audio encoders, and multi-face tracking.

\bibliographystyle{IEEEtran}
\bibliography{references}
\end{document}
